% !TeX encoding = UTF-8
\documentclass[institutional,screen,none]{ua-engineering-v6-beamer}
% !TeX encoding = UTF-8
\UASetUniversity{Universidad Austral}
\UASetFaculty{Facultad de Ingeniería}
\UASetDepartment{Departamento de Computación}
\UASetCampus{Pilar}
\UASetCareer{Ingeniería Informática}
\UASetCommission{Comisión A}
\UASetProfessor{Equipo docente}
\UASetSemester{Segundo cuatrimestre 2026}
\UASetCourse{Sistemas Distribuidos}
\UASetDate{2026}

\UASetSemester{Segundo cuatrimestre 2026}
\UASetDate{2026}
\usetikzlibrary{decorations.pathreplacing,matrix,backgrounds,shapes.multipart}

% Evitar cortes silábicos innecesarios y conservar una lectura fluida.
\hyphenpenalty=9000
\exhyphenpenalty=9000
\tolerance=1800
\emergencystretch=1.3em


% Etiqueta pedagógica visible. Se define localmente para no modificar el .sty
% compartido ni el build de GitHub.
\newcommand{\UASlideTag}[2]{%
  \begin{tikzpicture}[remember picture,overlay]
    \node[anchor=north east,rounded corners=1mm,inner xsep=3pt,inner ysep=1.5pt,
      draw=#2,fill=#2!8,text=#2,font=\fontsize{5.8}{6.4}\selectfont\bfseries]
      at ([xshift=-0.72cm,yshift=-1.08cm]current page.north east) {#1};
  \end{tikzpicture}%
}

\UASetClassNumber{06}
\UASetClassTitle{Storage distribuido y almacenamiento permanente}
\title{\UACourse}
\subtitle{Clase 06 --- Storage distribuido:\\almacenamiento permanente, acceso concurrente y rendimiento bajo fallas}
\author{\UAProfessor}
\date{\UADate}

\tikzset{
  uaNode/.style={draw=UAIngNavy,rounded corners=1.5mm,thick,fill=white,
    minimum height=0.72cm,align=center,font=\scriptsize,inner sep=3pt},
  uaSoft/.style={uaNode,fill=UAIngPaper},
  uaOrange/.style={uaNode,draw=UAIngOrange,fill=UAIngPaper},
  uaConcept/.style={uaNode,text width=3.15cm,minimum height=1.55cm,align=left},
  uaArrow/.style={->,very thick,>=Latex,draw=UAIngNavy},
  uaOrangeArrow/.style={->,very thick,>=Latex,draw=UAIngOrange},
  uaLabel/.style={font=\scriptsize,align=center,text=UAIngText},
  uaTiny/.style={font=\tiny,align=center,text=UAIngText},
  uaLane/.style={draw=UAIngRule,line width=0.9pt},
  uaEvent/.style={circle,fill=UAIngNavy,inner sep=2.2pt},
  uaEventOrange/.style={circle,fill=UAIngOrange,inner sep=2.4pt},
  uaEventBox/.style={draw=UAIngNavy,rounded corners=1.2mm,fill=white,
    minimum height=0.46cm,align=center,font=\tiny,inner xsep=3pt,inner ysep=1.5pt},
  uaEventBoxOrange/.style={uaEventBox,draw=UAIngOrange,fill=UAIngPaper},
  uaProcess/.style={draw=UAIngOrange,rounded corners=1.2mm,fill=UAIngPaper,
    minimum height=0.48cm,align=center,font=\tiny,inner xsep=4pt,inner ysep=1.7pt}
}

\newcommand{\SourceLine}[1]{\vfill{\tiny\color{UAIngMuted}Fuente / referencia: #1}}

\begin{document}
{
\setbeamertemplate{footline}{}
\begin{frame}[plain]
\titlepage
\end{frame}
}


% =============================================================================
\begin{frame}{Pregunta central: ¿qué resuelve el storage distribuido?}
\UASlideTag{NÚCLEO}{UAIngNavy}
\scriptsize
\begin{columns}[T]
\column{0.55\textwidth}
\begin{uakeybox}
En esta clase, \textbf{storage distribuido} significa \textbf{almacenamiento distribuido}: nodos y protocolos que convierten operaciones en estado permanente, recuperable y consultable aun cuando ocurren fallas.
\end{uakeybox}
\begin{itemize}
\item Conserva decisiones después de un reinicio.
\item Coordina lecturas y escrituras simultáneas.
\item Encuentra datos con latencia predecible.
\item Absorbe carga sin acumular trabajo sin límite.
\end{itemize}
\column{0.41\textwidth}
\begin{center}
\begin{tikzpicture}[node distance=0.23cm]
\node[uaOrange,text width=3.15cm,minimum height=0.70cm] (op) {Operación lógica\\``confirmar compra''};
\node[uaSoft,text width=3.15cm,below=of op] (dist) {Capa distribuida\\particiona, replica y ordena};
\node[uaSoft,text width=3.15cm,below=of dist] (eng) {Motores locales\\persisten, versionan e indexan};
\node[uaOrange,text width=3.15cm,below=of eng] (state) {Estado permanente\\recuperable y consultable};
\draw[uaOrangeArrow] (op)--(dist);
\draw[uaArrow] (dist)--(eng);
\draw[uaOrangeArrow] (eng)--(state);
\end{tikzpicture}
\end{center}
\end{columns}
\end{frame}

\begin{frame}{Dónde se ubica el storage dentro del sistema distribuido}
\UASlideTag{NÚCLEO}{UAIngNavy}
\scriptsize
\begin{center}
\begin{tikzpicture}[node distance=0.34cm and 0.42cm]
\node[uaOrange,text width=1.80cm] (clients) {Clientes};
\node[uaSoft,text width=2.05cm,right=of clients] (services) {Servicios de aplicación};
\node[uaSoft,text width=2.85cm,right=of services] (dist) {Capa distribuida\\enrutamiento, réplica y autoridad};
\node[uaSoft,text width=2.05cm,right=0.48cm of dist,yshift=1.2cm] (n1) {Nodo A\\motor + medio persistente};
\node[uaSoft,text width=2.05cm,right=0.48cm of dist] (n2) {Nodo B\\motor + medio persistente};
\node[uaSoft,text width=2.05cm,right=0.48cm of dist,yshift=-1.2cm] (n3) {Nodo C\\motor + medio persistente};
\draw[uaOrangeArrow] (clients)--(services);
\draw[uaArrow] (services)--(dist);
\draw[uaArrow] (dist.east)--(n1.west);
\draw[uaArrow] (dist.east)--(n2.west);
\draw[uaArrow] (dist.east)--(n3.west);
\node[draw=UAIngOrange,rounded corners=2mm,fit=(n1)(n2)(n3),inner sep=3pt,label={[uaTiny,text=UAIngOrange]above:nodos de almacenamiento}] {};
\end{tikzpicture}
\end{center}
\begin{columns}[T]
\column{0.48\textwidth}
\begin{block}{Capa distribuida}
Ubica y replica los datos; además, decide qué nodo tiene autoridad para confirmar cada operación.
\end{block}
\column{0.48\textwidth}
\begin{block}{Motor local}
Convierte esas decisiones en bytes persistentes, versiones visibles, índices y procedimientos de recuperación.
\end{block}
\end{columns}
\end{frame}

\begin{frame}{Qué llamaremos ``motor de almacenamiento''}
\UASlideTag{NÚCLEO}{UAIngNavy}
\begin{center}
\begin{tikzpicture}[node distance=0.45cm and 0.55cm]
\node[uaOrange,text width=2.0cm] (req) {Lectura o escritura};
\node[uaSoft,text width=2.35cm,right=of req] (txn) {Gestor transaccional};
\node[uaSoft,text width=2.35cm,right=of txn,yshift=1.05cm] (plan) {Planificador y ejecutor};
\node[uaSoft,text width=2.35cm,right=of txn] (buf) {Memoria intermedia};
\node[uaSoft,text width=2.35cm,right=of txn,yshift=-1.05cm] (log) {Registro de recuperación};
\node[uaSoft,text width=2.55cm,right=0.75cm of buf,yshift=0.65cm] (struct) {Árboles e índices\\archivos ordenados};
\node[uaOrange,text width=2.55cm,right=0.75cm of buf,yshift=-0.65cm] (media) {Páginas y archivos\\en medio persistente};
\draw[uaOrangeArrow] (req)--(txn);
\draw[uaArrow] (txn)--(plan);
\draw[uaArrow] (txn)--(buf);
\draw[uaArrow] (txn)--(log);
\draw[uaArrow] (plan)--(struct);
\draw[uaArrow] (buf)--(struct);
\draw[uaArrow] (log)--(media);
\draw[uaOrangeArrow] (struct)--(media);
\node[draw=UAIngNavy,rounded corners=2mm,fit=(txn)(plan)(buf)(log)(struct)(media),inner sep=6pt,label={[uaTiny]above:interior de un nodo de almacenamiento}] {};
\end{tikzpicture}
\end{center}
\begin{block}{Definición de trabajo}
El \textbf{motor de almacenamiento} es el software que traduce operaciones lógicas en bytes, versiones, estructuras de acceso y reglas de recuperación sobre un medio persistente.
\end{block}
\end{frame}

\begin{frame}{Siglas centrales I: WAL y MVCC}
\UASlideTag{NÚCLEO}{UAIngNavy}
\footnotesize
\begin{uainfo}
\textbf{WAL --- \textit{Write-Ahead Logging}.} En español: \textbf{registro anticipado de escritura}. El motor registra primero la evidencia necesaria para reconstruir un cambio después de una falla.
\end{uainfo}
\vspace{0.08cm}
\begin{uainfo}
\textbf{MVCC --- \textit{Multi-Version Concurrency Control}.} En español: \textbf{control de concurrencia multiversión}. El motor conserva varias versiones de un dato y decide cuál puede observar cada transacción.
\end{uainfo}
\vspace{0.08cm}
\begin{block}{Relación con el problema general}
WAL sostiene la recuperación de decisiones confirmadas. MVCC organiza el acceso concurrente a una historia de datos. Ninguno sustituye al otro.
\end{block}
\end{frame}

\begin{frame}{Siglas centrales II: LSM-tree}
\UASlideTag{NÚCLEO}{UAIngNavy}
\footnotesize
\begin{uainfo}
\textbf{LSM-tree --- \textit{Log-Structured Merge-tree}.} En español: \textbf{árbol de fusión estructurado como log}. Recibe escrituras en memoria, las vuelca a archivos ordenados e inmutables y, más tarde, fusiona esos archivos.
\end{uainfo}
\begin{columns}[T]
\column{0.48\textwidth}
\begin{block}{Qué optimiza}
Reduce el trabajo aleatorio en el camino crítico y favorece una alta tasa de ingestión.
\end{block}
\column{0.48\textwidth}
\begin{alertblock}{Qué costo desplaza}
Genera compactación posterior, lecturas sobre varios archivos y trabajo pendiente que debe controlarse.
\end{alertblock}
\end{columns}
\vspace{0.05cm}
\textbf{Convención.} Usaremos la sigla LSM, pero siempre indicaremos qué responsabilidad cumple, qué garantía aporta y qué costo desplaza.
\end{frame}

\begin{frame}{Cinco responsabilidades, no cinco términos aislados}
\UASlideTag{NÚCLEO}{UAIngNavy}
\scriptsize
\begin{center}
\begin{tabular}{p{0.25\textwidth}p{0.28\textwidth}p{0.34\textwidth}}
\toprule
Responsabilidad & Pregunta concreta & Mecanismo que estudiaremos \\
\midrule
Registrar decisiones & ¿Cómo sobrevive una compra ya confirmada? & Registro anticipado de escritura (WAL) y recuperación \\
Controlar concurrencia & ¿Qué puede ver un reporte mientras cambian los datos? & Control de concurrencia multiversión (MVCC) e instantáneas lógicas \\
Elegir una ruta física & ¿Cómo encontrar 20 filas entre un millón? & Índices, planificador y \texttt{EXPLAIN} (muestra el plan elegido) \\
Optimizar escrituras & ¿Cómo convertir muchas escrituras en entrada/salida eficiente? & Árbol B o árbol LSM \\
Proteger capacidad & ¿Qué ocurre si el mantenimiento queda atrás? & Compactación, trabajo pendiente y pausa de escritura \\
\bottomrule
\end{tabular}
\end{center}
\begin{uakeybox}
Ningún mecanismo constituye por sí solo el sistema de almacenamiento completo. Cada uno cubre una responsabilidad y aporta una garantía específica.
\end{uakeybox}
\end{frame}

\begin{frame}{Del dato lógico a los objetos físicos}
\UASlideTag{NÚCLEO}{UAIngNavy}
\begin{center}
\begin{tikzpicture}[node distance=0.32cm and 0.48cm]
\node[uaOrange,text width=2.45cm] (id) {Solicitud de compra\\ID 9F2A};
\node[uaSoft,text width=2.25cm,right=of id] (row) {Fila nueva en\\\texttt{purchases}};
\node[uaSoft,text width=2.25cm,right=of row] (page) {Página P17\\bloque físico};
\node[uaSoft,text width=2.25cm,right=of page] (media) {Archivo de datos\\en el medio};
\node[uaSoft,text width=2.25cm,below=0.50cm of row] (wal) {Entrada WAL\\para ID 9F2A};
\node[uaSoft,text width=2.25cm,below=0.50cm of page] (idx) {Entrada de índice I5\\ruta hacia P17};
\draw[uaOrangeArrow] (id)--(row);
\draw[uaArrow] (row)--(page);
\draw[uaArrow] (page)--(media);
\draw[uaArrow] (row)--(wal);
\draw[uaArrow] (idx)--(page);
\end{tikzpicture}
\end{center}
\begin{columns}[T]
\column{0.48\textwidth}
\begin{block}{Qué significa 9F2A}
Es un \textbf{identificador estable de la operación lógica}. Si el cliente reintenta, vuelve a enviar 9F2A para que el servidor reconozca la misma compra.
\end{block}
\column{0.48\textwidth}
\begin{block}{Qué no significa}
No es el contenido de la fila ni una posición física. P17 e I5 son etiquetas didácticas para una página de datos y una entrada de índice.
\end{block}
\end{columns}
\end{frame}

\begin{frame}{Vocabulario previo I: confirmar y volver durable}
\UASlideTag{NÚCLEO}{UAIngNavy}
\scriptsize
\begin{center}
\begin{tabular}{p{0.22\textwidth}p{0.65\textwidth}}
\toprule
Término & Significado en esta clase \\
\midrule
\textbf{COMMIT} & Decisión del motor de aceptar definitivamente los efectos de una transacción. El COMMIT puede estar decidido antes de que todas las páginas hayan sido escritas. \\
\textbf{ACK} & \textit{Acknowledgement}, o acuse de recibo. Es la respuesta enviada al cliente cuando se cumplió la garantía que el sistema promete. \\
\textbf{X durable} & La evidencia necesaria para reconstruir X sobrevivirá a la falla que el modelo declara. ``Durable'' siempre debe leerse respecto de una falla concreta. \\
\textbf{Log durable} & Porción del WAL cuyos bytes necesarios para recuperación ya sobrevivirían a la falla declarada. \\
\textbf{Flush durable} & Operación de volcado ya concluida que llevó los bytes requeridos del WAL más allá de la frontera durable. \\
\bottomrule
\end{tabular}
\end{center}
\begin{alertblock}{Regla de lectura}
Cuando aparezca ``X durable'', pregunte: ¿qué falla se considera y qué evidencia de X seguirá disponible después de esa falla?
\end{alertblock}
\end{frame}

\begin{frame}{Vocabulario previo II: escribir y vaciar buffers}
\UASlideTag{NÚCLEO}{UAIngNavy}
\scriptsize
\begin{center}
\begin{tabular}{p{0.22\textwidth}p{0.65\textwidth}}
\toprule
Término & Significado en esta clase \\
\midrule
\texttt{write}/append & Copiar o agregar bytes a una capa del sistema; la llamada puede terminar mientras los datos siguen en memoria o caché. \\
Flush o volcado & Mover datos acumulados en un buffer hacia una capa inferior del sistema de almacenamiento. Un flush común no implica necesariamente durabilidad. \\
\texttt{fsync} & Solicitar al sistema operativo que sincronice los datos relevantes con el medio persistente, de acuerdo con el hardware y la configuración. \\
Entrada/salida (E/S) & Transferencia de datos entre memoria, sistema operativo y dispositivo persistente. \\
Página sucia (dirty page) & Página modificada en memoria que todavía no fue escrita en el archivo de datos. \\
\bottomrule
\end{tabular}
\end{center}
\begin{alertblock}{Pregunta de ingeniería}
Cuando alguien diga ``ya está escrito'', pregunte: ¿en qué capa terminó la operación y qué falla todavía podría borrar esos bytes?
\end{alertblock}
\end{frame}

\begin{frame}{Vocabulario previo III: fallar y recuperar}
\UASlideTag{NÚCLEO}{UAIngNavy}
\scriptsize
\begin{center}
\begin{tabular}{p{0.22\textwidth}p{0.65\textwidth}}
\toprule
Término & Significado en esta clase \\
\midrule
Caída abrupta (crash) & Pérdida repentina de un proceso, del sistema operativo o de energía. El modelo debe indicar qué capas se pierden. \\
Recuperación (recovery) & Procedimiento de reinicio que usa la evidencia persistida para reconstruir un estado lógico válido. \\
REDO & Reaplicación de un cambio confirmado que aún no estaba materializado en su página de datos o índice. \\
Punto de control (checkpoint) & Punto durable que acota desde qué posición del WAL debe comenzar la recuperación. \\
LSN & \textit{Log Sequence Number}: número de secuencia que identifica una posición en el WAL; no es un reloj físico. \\
\bottomrule
\end{tabular}
\end{center}
\begin{uakeybox}
Estos términos describen cómo el motor transforma bytes persistidos en un estado lógico válido después de una falla.
\end{uakeybox}
\end{frame}

\begin{frame}{Vocabulario previo IV: concurrencia y consistencia}
\UASlideTag{NÚCLEO}{UAIngNavy}
\small
\begin{center}
\begin{tabular}{p{0.22\textwidth}p{0.64\textwidth}}
\toprule
Término & Significado en esta clase \\
\midrule
Versión & Representación de un dato creada por una transacción y visible para determinados lectores. \\
Snapshot & Instantánea lógica que determina qué transacciones y versiones puede observar una lectura. \\
Vista coherente & Resultado cuyas filas obedecen la misma regla de visibilidad, sin mezclar momentos arbitrarios. \\
Serializabilidad & Garantía de que el resultado concurrente equivale a alguna ejecución de las transacciones una detrás de otra. \\
\bottomrule
\end{tabular}
\end{center}
\end{frame}

\begin{frame}{Vocabulario previo V: consultas y definición del esquema}
\UASlideTag{NÚCLEO}{UAIngNavy}
\small
\begin{center}
\begin{tabular}{p{0.22\textwidth}p{0.64\textwidth}}
\toprule
Término & Significado en esta clase \\
\midrule
SQL & \textit{Structured Query Language}: lenguaje estructurado de consultas para leer y modificar datos. \\
DDL & \textit{Data Definition Language}: parte de SQL que crea o modifica estructuras, por ejemplo \texttt{CREATE INDEX}. \\
Planificador & Componente que compara rutas físicas posibles y elige un plan estimado. \\
\texttt{EXPLAIN} & Comando que muestra el plan elegido por el planificador. \\
\texttt{EXPLAIN ANALYZE} & Ejecuta la consulta y agrega mediciones reales al plan. \\
\bottomrule
\end{tabular}
\end{center}
\end{frame}

% =============================================================================
\UASectionSlide{Repaso ágil: cinco puentes\\hacia el almacenamiento}

\begin{frame}{Una frase por cada clase previa}
\UASlideTag{NÚCLEO}{UAIngNavy}
\scriptsize
\begin{center}
\begin{tikzpicture}[node distance=0.09cm]
\node[uaSoft,text width=0.91\textwidth,align=left] (a) {\textbf{Clase 1 --- Incertidumbre:} si no sabemos qué ocurrió, el almacenamiento conserva evidencia para responder de manera estable.};
\node[uaSoft,text width=0.91\textwidth,align=left,below=of a] (b) {\textbf{Clase 2 --- RPC y reintento:} el reintento es seguro solo si identidad, efecto y resultado sobreviven a la falla.};
\node[uaSoft,text width=0.91\textwidth,align=left,below=of b] (c) {\textbf{Clase 3 --- Replicación:} varias copias no indican por sí solas cuáles son durables, actuales ni autorizadas.};
\node[uaSoft,text width=0.91\textwidth,align=left,below=of c] (d) {\textbf{Clase 4 --- Consenso:} el protocolo decide qué debe recordarse; cada nodo debe persistirlo realmente.};
\node[uaSoft,text width=0.91\textwidth,align=left,below=of d] (e) {\textbf{Clase 5 --- Transacciones:} ACID define el contrato; WAL y MVCC implementan parte de ese contrato.};
\end{tikzpicture}
\end{center}
\vspace{0.05cm}
\textbf{Transición.} Hoy preguntamos qué bytes, esperas y estructuras vuelven verdaderas palabras como ``COMMIT'', ``durable'' y ``aislado''.
\end{frame}

\begin{frame}{Prediagnóstico intuitivo: razone antes de nombrar mecanismos}
\UASlideTag{ACTIVIDAD}{UAIngOrange}
\small
\begin{enumerate}
\item La aplicación mostró ``compra confirmada'' y un segundo después se cortó la energía. ¿Puede desaparecer al reiniciar?
\item Un reporte empezó a las 12:00 y durante su ejecución entraron nuevas compras. ¿Deberían cambiar sus totales a mitad del reporte?
\item Para mostrar 20 compras entre un millón, ¿es razonable revisar siempre todas las filas?
\item Ingresan 90 MB/s y el mantenimiento físico procesa 60 MB/s. ¿Puede la cola interna permanecer acotada indefinidamente?
\item Tres nodos conservan el dato solo en memoria principal (RAM). ¿Sobrevive un corte de energía común?
\end{enumerate}
\begin{alertblock}{Consigna}
Todavía no responda con el nombre de un mecanismo. Describa primero el comportamiento correcto y la falla que está imaginando.
\end{alertblock}
\end{frame}

\begin{frame}{Mapa conceptual: dónde aparece cada mecanismo}
\UASlideTag{NÚCLEO}{UAIngNavy}
\begin{center}
\begin{tikzpicture}[node distance=0.34cm and 0.38cm]
\node[uaOrange,text width=2.22cm,minimum height=1.25cm] (a) {Decisión durable\\\textbf{WAL + recuperación}};
\node[uaSoft,text width=2.22cm,minimum height=1.25cm,right=of a] (b) {Historia accedida\\concurrentemente\\\textbf{MVCC}};
\node[uaSoft,text width=2.22cm,minimum height=1.25cm,right=of b] (c) {Ruta de acceso\\\textbf{índice + EXPLAIN}};
\node[uaSoft,text width=2.22cm,minimum height=1.25cm,right=of c] (d) {Camino de escritura\\\textbf{árbol B / LSM}};
\node[uaOrange,text width=2.22cm,minimum height=1.25cm,right=of d] (e) {Protección de capacidad\\\textbf{compactación + pausa}};
\draw[uaOrangeArrow] (a)--(b);
\draw[uaArrow] (b)--(c);
\draw[uaArrow] (c)--(d);
\draw[uaOrangeArrow] (d)--(e);
\end{tikzpicture}
\end{center}
\begin{uakeybox}
El hilo conductor será siempre el mismo: problema, mecanismo, garantía, costo y evidencia.
\end{uakeybox}
\end{frame}

\begin{frame}{Evidencias de aprendizaje de esta clase}
\UASlideTag{NÚCLEO}{UAIngNavy}
\begin{columns}[T]
\column{0.48\textwidth}
\begin{block}{Durante la clase}
\begin{itemize}
\item completar un timeline WAL y decidir cuándo enviar el ACK;
\item completar un timeline MVCC y justificar qué versión es visible;
\item leer un plan \texttt{EXPLAIN} antes y después de un índice.
\end{itemize}
\end{block}
\column{0.48\textwidth}
\begin{block}{Al cerrar}
\begin{itemize}
\item diagnosticar backlog y write stall;
\item elegir un árbol B, una LSM-tree o una combinación según la carga de trabajo;
\item defender la decisión con falla, garantía, costo y métrica.
\end{itemize}
\end{block}
\end{columns}
\end{frame}

% =============================================================================
\UASectionSlide{WAL y recuperación:\\hacer durable una decisión}

\begin{frame}{CampusShop: qué cambia físicamente al confirmar una compra}
\UASlideTag{NÚCLEO}{UAIngNavy}
\scriptsize
\begin{columns}[T]
\column{0.54\textwidth}
\begin{uakeybox}
La operación con ID \textbf{9F2A} inserta una fila en \texttt{purchases}, descuenta una unidad en \texttt{inventory} y actualiza un índice por estudiante y fecha.
\end{uakeybox}
\begin{itemize}
\item 9F2A identifica la misma compra lógica si el cliente reintenta.
\item Las filas viven lógicamente en tablas.
\item El motor modifica páginas en memoria.
\item El WAL conserva evidencia para la recuperación.
\end{itemize}
\column{0.42\textwidth}
\begin{center}
\begin{tikzpicture}[node distance=0.22cm]
\node[uaOrange,text width=3.25cm] (op) {Solicitud\\ID de operación 9F2A};
\node[uaSoft,text width=3.25cm,below=of op] (p) {\texttt{purchases}: nueva fila};
\node[uaSoft,text width=3.25cm,below=of p] (i) {\texttt{inventory}: stock $-1$};
\node[uaSoft,text width=3.25cm,below=of i] (idx) {índice: nueva entrada};
\node[uaOrange,text width=3.25cm,below=of idx] (wal) {WAL: cambios + COMMIT};
\draw[uaArrow] (op)--(p);
\draw[uaArrow] (p)--(i);
\draw[uaArrow] (i)--(idx);
\draw[uaOrangeArrow] (idx)--(wal);
\end{tikzpicture}
\end{center}
\end{columns}
\end{frame}

\begin{frame}{La frontera durable: definición antes de usarla}
\UASlideTag{NÚCLEO}{UAIngNavy}
\small
\begin{columns}[T]
\column{0.42\textwidth}
\begin{center}
\begin{tikzpicture}[node distance=0.25cm]
\node[uaSoft,minimum width=3.0cm] (app) {Aplicación};
\node[uaSoft,below=of app,minimum width=3.0cm] (proc) {Buffer del proceso};
\node[uaSoft,below=of proc,minimum width=3.0cm] (os) {Caché del sistema operativo};
\node[uaSoft,below=of os,minimum width=3.0cm] (dev) {Caché del dispositivo};
\node[uaOrange,below=0.55cm of dev,minimum width=3.0cm] (media) {Medio persistente};
\draw[uaArrow] (app)--(proc);
\draw[uaArrow] (proc)--(os);
\draw[uaArrow] (os)--(dev);
\draw[uaOrangeArrow] (dev)--node[right,uaTiny,text=UAIngOrange]{cruce durable}(media);
\end{tikzpicture}
\end{center}
\column{0.54\textwidth}
\begin{block}{Definición}
La \textbf{frontera durable} es el punto a partir del cual, bajo el modelo de falla declarado, permanece disponible la evidencia necesaria para reconstruir la operación.
\end{block}
\begin{alertblock}{No existe una frontera universal}
Depende de la falla considerada y de la semántica efectiva de \texttt{fsync}, del sistema operativo, del controlador, del dispositivo y de su configuración.
\end{alertblock}
\end{columns}
\end{frame}

\begin{frame}{Qué es WAL y por qué está separado de las páginas}
\UASlideTag{NÚCLEO}{UAIngNavy}
\scriptsize
\begin{columns}[T]
\column{0.47\textwidth}
\begin{block}{Definición intuitiva}
WAL es el registro secuencial del motor. Primero conserva la evidencia que permitirá reconstruir una decisión; después, las páginas pueden escribirse en el orden físico más conveniente.
\end{block}
\begin{block}{Definición técnica}
El registro anticipado de escritura impone dos órdenes: la evidencia de recuperación debe volverse durable antes que la página modificada; además, el COMMIT debe volverse durable antes de enviar el ACK.
\end{block}
\column{0.49\textwidth}
\begin{center}
\begin{tikzpicture}[x=0.92cm,y=0.92cm]
\node[uaOrange,text width=3.35cm] (txn) at (0,2.35) {Gestor transaccional};
\node[uaSoft,text width=3.05cm] (wal1) at (-2.05,1.10) {1. Agrega cambios\\al archivo WAL};
\node[uaSoft,text width=3.05cm] (buf) at (2.05,1.10) {2. Modifica páginas\\en memoria};
\node[uaOrange,text width=3.05cm] (wald) at (-2.05,-0.20) {Log durable\\después del flush};
\node[uaSoft,text width=3.05cm] (pages) at (2.05,-0.20) {Páginas persistidas\\más adelante};
\node[uaSoft,text width=4.0cm] (rec) at (0,-1.55) {Recuperación: usa el WAL durable\\para completar páginas atrasadas};
\draw[uaOrangeArrow] (txn.south west)--(wal1.north);
\draw[uaArrow] (txn.south east)--(buf.north);
\draw[uaOrangeArrow] (wal1)--(wald);
\draw[uaArrow] (buf)--(pages);
\draw[uaArrow] (wald.south)--(rec.north west);
\draw[uaOrangeArrow] (rec.north east)--(pages.south);
\end{tikzpicture}
\end{center}
\end{columns}
\end{frame}

\begin{frame}{WAL: dos reglas, dos fallas que evitan}
\UASlideTag{NÚCLEO}{UAIngNavy}
\scriptsize
\begin{columns}[T]
\column{0.48\textwidth}
\begin{block}{Regla 1 --- log antes que página}
\textbf{Orden:} registro WAL $\rightarrow$ flush durable $\rightarrow$ página de datos.
\end{block}
\begin{itemize}
\item \textbf{Propósito:} que toda página persistida tenga evidencia suficiente para ser interpretada o rehecha.
\item \textbf{Falla evitada:} una página nueva sin registro recuperable del cambio que la produjo.
\end{itemize}
\column{0.48\textwidth}
\begin{block}{Regla 2 --- COMMIT antes que ACK}
\textbf{Orden:} registro COMMIT $\rightarrow$ flush durable $\rightarrow$ ACK al cliente.
\end{block}
\begin{itemize}
\item \textbf{Propósito:} que cada confirmación enviada pueda reconstruirse tras un crash.
\item \textbf{Falla evitada:} mostrar ``confirmada'' y perder luego la decisión de COMMIT.
\end{itemize}
\end{columns}
\begin{uakeybox}
Garantía buscada: después de la recuperación aparecen todas las transacciones confirmadas y ninguna transacción incompleta se presenta como confirmada.
\end{uakeybox}
\end{frame}

\begin{frame}{Por qué WAL aparece como un carril propio en el timeline}
\UASlideTag{NÚCLEO}{UAIngNavy}
\small
\begin{center}
\begin{tikzpicture}[x=1cm,y=1cm]
\node[uaOrange,text width=1.65cm] (client) at (-5.85,0.35) {Cliente};
\node[uaSoft,text width=1.95cm] (service) at (-3.15,0.35) {Servicio de aplicación};
\node[draw=UAIngOrange,rounded corners=2mm,minimum width=8.05cm,minimum height=2.85cm] at (1.80,0.20) {};
\node[uaTiny,text=UAIngOrange,anchor=south] at (1.80,1.72) {un único motor de almacenamiento};
\node[uaSoft,text width=2.05cm] (txn) at (-0.65,1.00) {Gestor transaccional};
\node[uaSoft,text width=2.05cm] (wal) at (2.05,1.00) {Archivo WAL};
\node[uaSoft,text width=2.05cm] (mem) at (-0.65,-0.55) {Páginas en memoria};
\node[uaOrange,text width=2.20cm] (pages) at (2.05,-0.55) {Páginas e índices persistentes};
\node[uaSoft,text width=1.85cm] (rec) at (4.80,0.20) {Recuperación};
\draw[uaOrangeArrow] (client)--(service);
\draw[uaArrow] (service.east)--(txn.west);
\draw[uaArrow] (txn)--(wal);
\draw[uaArrow] (txn)--(mem);
\draw[uaOrangeArrow] (mem)--(pages);
\draw[uaArrow] (wal.east)--(rec.north west);
\draw[uaOrangeArrow] (rec.south west)--(pages.east);
\end{tikzpicture}
\end{center}
\vspace{-0.10cm}
\begin{columns}[T]
\column{0.48\textwidth}
\begin{itemize}
\item \textbf{1. Registrar:} el gestor agrega evidencia al WAL.
\item \textbf{2. Modificar:} las páginas cambian primero en memoria.
\end{itemize}
\column{0.48\textwidth}
\begin{itemize}
\item \textbf{3. Materializar:} las páginas se escriben después.
\item \textbf{4. Recuperar:} tras un crash, el WAL permite completarlas.
\end{itemize}
\end{columns}

\end{frame}

\begin{frame}{Timeline WAL vacío: ubique primero las fronteras}
\UASlideTag{ACTIVIDAD}{UAIngOrange}
\begin{center}
\begin{tikzpicture}[x=0.92cm,y=0.67cm]
\foreach \y/\lab in {4/{Cliente},3/{Servicio},2/{Motor / memoria},1/{Archivo WAL},0/{Páginas / índices}}{
  \node[uaTiny,anchor=east] at (0.55,\y){\lab};
  \draw[uaLane] (0.75,\y)--(10.6,\y);
}
\draw[->,thick,UAIngNavy] (0.75,-0.62)--(10.6,-0.62) node[right,uaTiny]{tiempo};
\node[uaLabel,text width=9.1cm] at (5.55,4.55){Ubique: solicitud, cambios, página sucia, registro COMMIT, flush durable, ACK, crash y REDO.};
\end{tikzpicture}
\end{center}
\begin{alertblock}{Predicción}
¿Qué evento debe haberse completado antes de mostrar ``compra confirmada''? Marque también qué carriles pertenecen al mismo motor de almacenamiento.
\end{alertblock}
\end{frame}

\begin{frame}{Construcción progresiva del timeline WAL}
\UASlideTag{NÚCLEO}{UAIngNavy}
\begin{center}
\begin{tikzpicture}[x=0.92cm,y=0.66cm]
\foreach \y/\lab in {4/{Cliente},3/{Servicio},2/{Motor / memoria},1/{WAL},0/{Páginas / índices}}{
  \node[uaTiny,anchor=east] at (0.55,\y){\lab};
  \draw[uaLane] (0.75,\y)--(10.6,\y);
}
\draw[->,thick,UAIngNavy] (0.75,-0.62)--(10.6,-0.62) node[right,uaTiny]{tiempo};

\uncover<1->{
  \node[uaEventBoxOrange] (e1) at (1.35,4) {1. solicitud\,9F2A};
  \draw[uaOrangeArrow] (e1.south)--(1.35,3.08);
}
\uncover<2->{
  \node[uaEventBox] (e2) at (2.35,2) {2. preparar\\cambios};
  \draw[uaArrow] (2.35,2.92)--(e2.north);
}
\uncover<3->{
  \node[uaEventBox] (e3) at (3.40,1) {3. WAL\\UPDATE};
  \draw[uaArrow] (3.40,1.92)--(e3.north);
}
\uncover<4->{
  \node[uaEventBox] (e4) at (4.45,0) {4. páginas\\sucias};
  \draw[uaArrow] (4.45,1.92)--(e4.north);
}
\uncover<5->{
  \node[uaEventBox] (e5) at (5.45,1) {5. registro\\COMMIT};
  \draw[uaArrow] (5.45,1.92)--(e5.north);
}
\uncover<6->{
  \draw[very thick,UAIngOrange] (6.55,-0.15)--(6.55,4.18);
  \node[uaEventBoxOrange,anchor=south] at (6.55,4.25) {6. flush durable};
}
\uncover<7->{
  \node[uaEventBoxOrange] (e7) at (7.55,4) {7. ACK\\confirmada};
  \draw[uaOrangeArrow] (7.55,3.08)--(e7.south);
}
\uncover<8->{
  \draw[very thick,UAIngNavy,dashed] (8.55,-0.15)--(8.55,4.18);
  \node[uaEventBox,anchor=south] at (8.55,4.25) {8. crash};
}
\uncover<9->{
  \node[uaEventBoxOrange] (e9) at (9.65,0) {9. REDO};
  \draw[uaOrangeArrow] (9.65,0.92)--(e9.north);
}
\end{tikzpicture}
\end{center}
\only<1>{\begin{block}{Paso 1}El cliente envía una intención con identidad estable.\end{block}}
\only<2>{\begin{block}{Paso 2}El motor prepara cambios lógicos en memoria.\end{block}}
\only<3>{\begin{block}{Paso 3}El WAL describe cómo rehacerlos; todavía puede estar en buffers.\end{block}}
\only<4>{\begin{block}{Paso 4}Las páginas se modifican en memoria y quedan sucias.\end{block}}
\only<5>{\begin{block}{Paso 5}El registro COMMIT expresa la decisión, pero aún no es durable.\end{block}}
\only<6>{\begin{block}{Paso 6}El flush lleva el WAL relevante más allá de la frontera durable declarada.\end{block}}
\only<7>{\begin{block}{Paso 7}Ahora el servicio puede confirmar sin esperar que todas las páginas estén escritas.\end{block}}
\only<8>{\begin{block}{Paso 8}El crash borra memoria y buffers volátiles.\end{block}}
\only<9>{\begin{block}{Paso 9}Recuperación lee el WAL y REDO materializa en las páginas los cambios confirmados que faltan.\end{block}}
\end{frame}

\begin{frame}{Criterio de corrección de recuperación: todas y solo las confirmadas}
\UASlideTag{NÚCLEO}{UAIngNavy}
\scriptsize
\begin{center}
\begin{tabular}{p{0.24\textwidth}p{0.28\textwidth}p{0.39\textwidth}}
\toprule
Punto del crash & Evidencia durable disponible & Resultado correcto después de recuperar \\
\midrule
Antes del COMMIT durable & Puede haber registros de cambio, pero no una decisión durable de COMMIT & La compra no aparece confirmada; el trabajo incompleto se descarta o revierte según el diseño. \\
Después del COMMIT durable, antes de escribir páginas & Registros de cambio y COMMIT en el WAL durable & REDO reconstruye pedido, stock e índice. \\
Después del COMMIT y de las páginas & WAL y páginas ya están materializados & La recuperación verifica el estado; repetir REDO no debe producir un segundo efecto lógico. \\
\bottomrule
\end{tabular}
\end{center}
\begin{uakeybox}
Los tres resultados son correctos porque aplican el mismo criterio: conservar todas las transacciones confirmadas y no convertir una transacción incompleta en confirmada.
\end{uakeybox}
\end{frame}

\begin{frame}{Group commit: compartir el costo sin debilitar la garantía}
\UASlideTag{NÚCLEO}{UAIngNavy}
\small
\begin{center}
\begin{tikzpicture}[node distance=0.25cm and 0.55cm]
\node[uaSoft,text width=2.15cm] (t1) {$T_1$ compra};
\node[uaSoft,text width=2.15cm,below=of t1] (t2) {$T_2$ pago};
\node[uaSoft,text width=2.15cm,below=of t2] (t3) {$T_3$ reserva};
\node[uaSoft,text width=2.85cm,right=0.75cm of t2] (buf) {Buffer WAL\\tres COMMIT};
\node[uaOrange,text width=2.35cm,right=0.55cm of buf] (fs) {un flush durable\\compartido};
\node[uaSoft,text width=2.35cm,right=0.55cm of fs] (acks) {tres ACK\\después del flush};
\draw[uaArrow] (t1.east)--(buf.west);
\draw[uaArrow] (t2.east)--(buf.west);
\draw[uaArrow] (t3.east)--(buf.west);
\draw[uaOrangeArrow] (buf)--(fs);
\draw[uaArrow] (fs)--(acks);
\end{tikzpicture}
\end{center}
\begin{uakeybox}
\textbf{Ganancia:} un acceso físico amortiza varios COMMIT. \quad
\textbf{Costo:} cada transacción puede esperar brevemente.\\[2pt]
\textbf{Garantía:} cada ACK sigue saliendo solamente después de que su COMMIT quedó durable.
\end{uakeybox}
\end{frame}

\begin{frame}{Checkpoint: un punto durable desde donde comenzar recuperación}
\UASlideTag{NÚCLEO}{UAIngNavy}
\small
\begin{center}
\begin{tikzpicture}[x=0.93cm,y=0.72cm]
\draw[->,thick,UAIngNavy] (0,0)--(9.4,0) node[right,uaTiny]{LSN / progreso};
\fill[UAIngPaper] (0.7,-0.42) rectangle (4.0,0.42);
\fill[UAIngOrange!18] (4.0,-0.42) rectangle (8.4,0.42);
\draw[very thick,UAIngNavy] (4.0,-0.65)--(4.0,0.80);
\draw[very thick,UAIngOrange] (8.4,-0.65)--(8.4,0.80);
\node[uaTiny] at (2.35,0){páginas materializadas};
\node[uaTiny,text=UAIngOrange] at (6.20,0){WAL que debe reproducirse};
\node[uaTiny,above=0.08cm] at (4.0,0.80){checkpoint};
\node[uaTiny,text=UAIngOrange,above=0.08cm] at (8.4,0.80){crash};
\end{tikzpicture}
\end{center}
\begin{block}{Qué es}
Registra suficiente progreso durable para que la recuperación no comience desde el origen del WAL.
\end{block}
\begin{itemize}
\item \textbf{Más frecuente:} menos WAL por reproducir, pero más entrada/salida y posibles picos.
\item \textbf{Más espaciado:} operación normal más suave, pero recuperación y retención mayores.
\end{itemize}
\end{frame}

\begin{frame}{Caso real: PostgreSQL desacopla COMMIT de las páginas}
\UASlideTag{NÚCLEO}{UAIngNavy}
\begin{columns}[T]
\column{0.58\textwidth}
\begin{itemize}
\item El WAL relevante debe llegar a almacenamiento estable antes que los cambios de archivos de datos.
\item Tras un crash, recuperación reproduce WAL para restaurar un estado válido.
\item Las páginas de tablas e índices pueden estar atrasadas respecto del COMMIT.
\item Hardware y configuración forman parte de la garantía efectiva.
\end{itemize}
\column{0.38\textwidth}
\begin{center}
\begin{tikzpicture}[node distance=0.35cm]
\node[uaSoft,minimum width=3.2cm] (wal) {WAL secuencial};
\node[uaOrange,below=of wal,minimum width=3.2cm] (flush) {COMMIT durable};
\node[uaSoft,below=of flush,minimum width=3.2cm] (pages) {páginas atrasadas};
\node[uaSoft,below=of pages,minimum width=3.2cm] (redo) {recuperación / REDO};
\draw[uaArrow] (wal)--(flush);
\draw[uaArrow] (flush)--(pages);
\draw[uaOrangeArrow] (pages)--(redo);
\end{tikzpicture}
\end{center}
\end{columns}
\SourceLine{PostgreSQL, documentación oficial sobre Write-Ahead Logging y recuperación.}
\end{frame}

\begin{frame}{Checkpoint de comprensión: WAL y recuperación}
\UASlideTag{ACTIVIDAD}{UAIngOrange}
\begin{enumerate}
\item Defina con sus palabras la frontera durable para el caso CampusShop.
\item Explique por qué una página puede seguir en memoria después de un COMMIT correcto.
\item Mueva el crash a tres puntos del timeline y prediga el estado recuperado.
\item Explique qué garantiza WAL y qué queda a cargo de replicación, backups e idempotencia.
\end{enumerate}
\end{frame}

% =============================================================================
\UASectionSlide{MVCC: acceder concurrentemente\\a una historia coherente}

\begin{frame}{Problema antes del mecanismo: un reporte y 900 compras nuevas}
\UASlideTag{NÚCLEO}{UAIngNavy}
\begin{uakeybox}
Un reporte del comedor comienza a las 12:00 y tarda 40 segundos. Durante ese intervalo entran 900 compras. ¿Debe mezclar filas de momentos distintos? ¿Debe bloquear a todos los escritores?
\end{uakeybox}
\begin{center}
\begin{tikzpicture}[node distance=0.55cm and 0.7cm]
\node[uaOrange,text width=2.8cm] (report) {Reporte largo\\inicia 12:00};
\node[uaSoft,text width=2.8cm,right=of report] (writes) {900 escrituras\\12:00--12:00:40};
\node[uaSoft,text width=3.1cm,right=of writes] (question) {¿Bloquear escritores, mezclar momentos o conservar una vista coherente?};
\draw[uaOrangeArrow] (report)--(writes);
\draw[uaArrow] (writes)--(question);
\end{tikzpicture}
\end{center}
\end{frame}

\begin{frame}{Vocabulario de MVCC antes del timeline}
\UASlideTag{NÚCLEO}{UAIngNavy}
\scriptsize
\begin{center}
\begin{tabular}{p{0.20\textwidth}p{0.68\textwidth}}
\toprule
Término & Definición de trabajo \\
\midrule
Versión & Representación de una fila creada por una transacción y válida para determinados lectores. \\
Transacción lectora & Operación que consulta datos bajo una regla de visibilidad. \\
Transacción escritora & Operación que crea una versión nueva y no borra de inmediato la anterior. \\
Snapshot & Instantánea lógica que determina qué transacciones confirmadas y qué versiones son visibles. \\
Visibilidad & Decisión de si una versión pertenece a la historia que ese snapshot puede observar. \\
Vista coherente & Resultado cuyas filas se seleccionan con la misma regla de visibilidad. \\
Serializabilidad & Garantía de que la ejecución concurrente equivale a alguna ejecución serial de las transacciones. \\
\bottomrule
\end{tabular}
\end{center}
\begin{alertblock}{Límite de la metáfora}
``Fotografía'' no significa copiar toda la base. Tampoco implica serializabilidad por sí sola: describe una regla para elegir versiones.
\end{alertblock}
\end{frame}

\begin{frame}{MVCC en una imagen: dos lectores, dos vistas coherentes}
\UASlideTag{NÚCLEO}{UAIngNavy}
\begin{center}
\begin{tikzpicture}[node distance=0.38cm and 0.55cm]
\node[uaOrange,text width=3.2cm] (before) {Antes\\v1 = 120 ya está confirmada};
\node[uaSoft,text width=3.4cm,right=of before] (writer) {Escritor\\crea v2 = 135 y ejecuta COMMIT};
\node[uaSoft,text width=3.25cm,right=of writer,yshift=0.65cm] (old) {Reporte iniciado antes\\sigue leyendo v1};
\node[uaSoft,text width=3.25cm,right=of writer,yshift=-0.65cm] (new) {Lectura iniciada después\\puede leer v2};
\draw[uaOrangeArrow] (before)--(writer);
\draw[uaArrow] (writer.east)--++(0.28,0)--(old.west);
\draw[uaOrangeArrow] (writer.east)--++(0.28,0)--(new.west);
\end{tikzpicture}
\end{center}
\begin{block}{Idea central}
MVCC conserva versiones y aplica reglas de visibilidad. Dos lectores pueden observar versiones diferentes sin que el reporte antiguo cambie a mitad de ejecución.
\end{block}
\end{frame}

\begin{frame}{Timeline MVCC vacío: cuatro carriles y términos conocidos}
\UASlideTag{ACTIVIDAD}{UAIngOrange}
\begin{center}
\begin{tikzpicture}[x=1.0cm,y=0.78cm]
\foreach \y/\lab in {3/{Lector $T_R$},2/{Escritor $T_W$},1/{Versiones},0/{Lectura visible}}{
  \node[uaTiny,anchor=east] at (0.65,\y){\lab};
  \draw[uaLane] (0.85,\y)--(10.1,\y);
}
\draw[->,thick,UAIngNavy] (0.85,-0.58)--(10.1,-0.58) node[right,uaTiny]{tiempo};
\node[uaLabel,text width=8.9cm] at (5.4,3.62){Ubique: v1, inicio del snapshot, lectura, creación de v2, COMMIT, segunda lectura, fin y reclamación.};
\end{tikzpicture}
\end{center}
\begin{alertblock}{Predicción}
Después del COMMIT de v2, ¿qué lee una segunda consulta dentro del mismo reporte y qué lee un reporte nuevo?
\end{alertblock}
\end{frame}

\begin{frame}{Construcción progresiva del timeline MVCC}
\UASlideTag{NÚCLEO}{UAIngNavy}
\begin{center}
\begin{tikzpicture}[x=1.0cm,y=0.76cm]
\foreach \y/\lab in {3/{Lector $T_R$},2/{Escritor $T_W$},1/{Versiones},0/{Lectura visible}}{
  \node[uaTiny,anchor=east] at (0.65,\y){\lab};
  \draw[uaLane] (0.85,\y)--(10.1,\y);
}
\draw[->,thick,UAIngNavy] (0.85,-0.60)--(10.1,-0.60) node[right,uaTiny]{tiempo};
\uncover<1->{\node[uaEventBox] at (1.35,1) {1. v1};}
\uncover<2->{
  \node[uaEventBoxOrange,text width=1.55cm] at (2.45,3) {2. inicia $T_R$\\snapshot $S_R$};
  \draw[very thick,UAIngOrange] (3.35,2.66)--(8.05,2.66);
}
\uncover<3->{
  \node[uaEventBox] (m3) at (3.55,0) {3. lee v1};
  \draw[uaArrow] (3.55,2.90)--(m3.north);
}
\uncover<4->{
  \node[uaEventBox] (m4a) at (4.65,2) {4. crea};
  \node[uaEventBox] (m4b) at (4.65,1) {v2};
  \draw[uaArrow] (m4a.south)--(m4b.north);
}
\uncover<5->{\node[uaEventBoxOrange] at (6.20,2) {5. COMMIT};}
\uncover<6->{
  \node[uaEventBox] at (6.95,0) {$T_R$: v1};
  \node[uaEventBoxOrange] at (6.95,1) {lector nuevo: v2};
}
\uncover<7->{\node[uaEventBoxOrange] at (8.30,3) {7. termina $T_R$};}
\uncover<8->{\node[uaEventBoxOrange] at (9.35,1) {8. v1 reclamable};}
\end{tikzpicture}
\end{center}
\only<1>{\begin{block}{Paso 1}v1 es la versión confirmada existente.\end{block}}
\only<2>{\begin{block}{Paso 2}$T_R$ fija una regla de visibilidad: su snapshot.\end{block}}
\only<3>{\begin{block}{Paso 3}La primera lectura elige v1.\end{block}}
\only<4>{\begin{block}{Paso 4}$T_W$ crea v2 sin borrar v1.\end{block}}
\only<5>{\begin{block}{Paso 5}v2 confirma y queda disponible para snapshots posteriores.\end{block}}
\only<6>{\begin{block}{Paso 6}$T_R$ mantiene v1; un lector nuevo puede elegir v2.\end{block}}
\only<7>{\begin{block}{Paso 7}Termina el snapshot largo.\end{block}}
\only<8>{\begin{block}{Paso 8}v1 puede limpiarse cuando ninguna regla de visibilidad la necesite.\end{block}}
\end{frame}

\begin{frame}{Regla simplificada de visibilidad}
\UASlideTag{NÚCLEO}{UAIngNavy}
\begin{center}
\begin{tikzpicture}[x=1.0cm,y=0.8cm]
\draw[->,thick,UAIngNavy] (0,0)--(9.1,0) node[right,uaTiny]{historia lógica};
\draw[very thick,UAIngNavy] (0.8,0.8)--(4.7,0.8);
\draw[very thick,UAIngOrange] (4.7,0.8)--(8.3,0.8);
\node[uaTiny,above=0.12cm] at (2.75,0.8){v1 visible};
\node[uaTiny,text=UAIngOrange,above=0.12cm] at (6.5,0.8){v2 visible};
\draw[UAIngNavy,dashed] (3.1,-0.25)--(3.1,1.4);
\draw[UAIngOrange,dashed] (6.3,-0.25)--(6.3,1.4);
\node[uaSoft,text width=2.15cm] at (3.1,-0.90){snapshot A\\elige v1};
\node[uaOrange,text width=2.15cm] at (6.3,-0.90){snapshot B\\elige v2};
\end{tikzpicture}
\end{center}
\begin{block}{Lectura conceptual}
Una versión es visible si su creador pertenece a la historia confirmada que el snapshot puede observar y su reemplazo/eliminación todavía no era visible para ese snapshot.
\end{block}
\end{frame}

\begin{frame}{Caso real: PostgreSQL, versiones de fila y VACUUM}
\UASlideTag{NÚCLEO}{UAIngNavy}
\begin{columns}[T]
\column{0.58\textwidth}
\begin{itemize}
\item Un \texttt{UPDATE} crea una versión nueva de la fila.
\item La versión anterior permanece mientras algún snapshot pueda verla.
\item \texttt{VACUUM} recupera espacio reutilizable cuando la historia deja de ser necesaria.
\item Una transacción larga retrasa la limpieza y aumenta bloat, E/S y presión sobre índices.
\end{itemize}
\column{0.38\textwidth}
\begin{center}
\begin{tikzpicture}[node distance=0.35cm]
\node[uaSoft,minimum width=3.25cm] (v1) {v1 visible};
\node[uaOrange,below=of v1,minimum width=3.25cm] (v2) {UPDATE crea v2};
\node[uaSoft,below=of v2,minimum width=3.25cm] (hold) {snapshot retiene v1};
\node[uaSoft,below=of hold,minimum width=3.25cm] (vac) {VACUUM cuando es seguro};
\draw[uaArrow] (v1)--(v2);
\draw[uaArrow] (v2)--(hold);
\draw[uaOrangeArrow] (hold)--(vac);
\end{tikzpicture}
\end{center}
\end{columns}
\SourceLine{PostgreSQL, documentación oficial sobre MVCC y routine vacuuming.}
\end{frame}

\begin{frame}{MVCC no equivale automáticamente a serializabilidad}
\UASlideTag{NÚCLEO}{UAIngNavy}
\begin{center}
\begin{tikzpicture}[node distance=0.45cm and 1.0cm]
\node[uaSoft,text width=3.5cm] (t1a) {$T_1$: lee A=1, B=1};
\node[uaSoft,text width=3.5cm,right=of t1a] (t2a) {$T_2$: lee A=1, B=1};
\node[uaSoft,text width=3.5cm,below=of t1a] (t1b) {$T_1$: escribe A=0};
\node[uaSoft,text width=3.5cm,below=of t2a] (t2b) {$T_2$: escribe B=0};
\node[uaOrange,text width=5.5cm,below=0.65cm of $(t1b)!0.5!(t2b)$] (end) {Ambas confirman: A=0, B=0\\se viola ``al menos una debe seguir activa''};
\draw[uaArrow] (t1a)--(t1b);
\draw[uaArrow] (t2a)--(t2b);
\draw[uaOrangeArrow] (t1b)--(end);
\draw[uaOrangeArrow] (t2b)--(end);
\end{tikzpicture}
\end{center}
\begin{block}{Conclusión}
MVCC administra versiones. El nivel de aislamiento y la detección de conflictos determinan qué anomalías, como write skew, siguen siendo posibles.
\end{block}
\end{frame}

\begin{frame}{Checkpoint de comprensión: MVCC}
\UASlideTag{ACTIVIDAD}{UAIngOrange}
\begin{enumerate}
\item Explique snapshot sin usar la metáfora de una copia completa.
\item ¿Qué lee un reporte que comenzó antes del COMMIT de v2?
\item ¿Cuándo puede limpiarse v1 y quién puede retrasar esa limpieza?
\item ¿Qué no se puede concluir solo porque el motor ``usa MVCC''?
\end{enumerate}
\end{frame}

% =============================================================================
\UASectionSlide{Índices y EXPLAIN:\\medir antes de optimizar}

\begin{frame}{Índice: una estructura adicional que evita leer todo}
\UASlideTag{NÚCLEO}{UAIngNavy}
\begin{columns}[T]
\column{0.42\textwidth}
\begin{center}
\begin{tikzpicture}[node distance=0.35cm]
\node[uaOrange,minimum width=3.6cm] (book) {Tabla: 1.000.000 compras};
\node[uaSoft,below=of book,minimum width=3.6cm] (idx) {Índice por estudiante + fecha};
\node[uaSoft,below=of idx,minimum width=3.6cm] (leaf) {Hoja: estudiante 42\\fechas descendentes};
\draw[uaArrow] (book)--(idx);
\draw[uaOrangeArrow] (idx)--(leaf);
\end{tikzpicture}
\end{center}
\column{0.54\textwidth}
\begin{block}{Intuición}
Como el índice de un libro, agrega una ruta corta hacia ciertos datos. No reemplaza el libro: ocupa espacio y debe actualizarse cuando cambia el contenido.
\end{block}
\begin{alertblock}{Consecuencia}
Acelera algunas preguntas y encarece cada escritura que deba mantenerlo.
\end{alertblock}
\end{columns}
\end{frame}

\begin{frame}{La consulta del caso conductor}
\UASlideTag{NÚCLEO}{UAIngNavy}
\begin{block}{Pregunta de producto}
Mostrar las últimas 20 compras confirmadas de una estudiante.
\end{block}
\begin{center}
\begin{tikzpicture}[node distance=0.55cm]
\node[uaSoft,text width=9.8cm,align=left,font=\scriptsize\ttfamily] (q) {
SELECT id, total, confirmed\_at\\
FROM purchases\\
WHERE student\_id = 8421 AND status = 'CONFIRMED'\\
ORDER BY confirmed\_at DESC\\
LIMIT 20;};
\node[uaOrange,text width=2.25cm,below=0.55cm of q.south west,anchor=north west] (a) {igualdad\\student\_id};
\node[uaSoft,text width=2.25cm,right=of a] (b) {igualdad\\status};
\node[uaSoft,text width=2.25cm,right=of b] (c) {orden\\fecha DESC};
\node[uaSoft,text width=2.25cm,right=of c] (d) {corte temprano\\LIMIT 20};
\draw[uaOrangeArrow] (a)--(b);
\draw[uaArrow] (b)--(c);
\draw[uaArrow] (c)--(d);
\end{tikzpicture}
\end{center}
\end{frame}

\begin{frame}{EXPLAIN y EXPLAIN ANALYZE: plan frente a evidencia}
\UASlideTag{NÚCLEO}{UAIngNavy}
\begin{columns}[T]
\column{0.47\textwidth}
\begin{block}{\texttt{EXPLAIN}}
Muestra el plan elegido y sus \textbf{estimaciones}: forma de acceso, filtros, orden, costo y cantidad esperada de filas.
\end{block}
\column{0.47\textwidth}
\begin{alertblock}{\texttt{EXPLAIN ANALYZE}}
Ejecuta la consulta y agrega \textbf{mediciones reales}: tiempos, filas, repeticiones y, con \texttt{BUFFERS}, entrada/salida observada.
\end{alertblock}
\end{columns}
\vspace{0.3cm}
\begin{center}
\begin{tikzpicture}[node distance=0.55cm]
\node[uaSoft,text width=2.5cm] (sql) {Consulta SQL};
\node[uaSoft,text width=2.5cm,right=of sql] (planificador) {Planificador\\elige una ruta};
\node[uaSoft,text width=2.5cm,right=of planificador] (plan) {Plan estimado};
\node[uaOrange,text width=2.5cm,right=of plan] (run) {Ejecución medida};
\draw[uaArrow] (sql)--(planificador);
\draw[uaArrow] (planificador)--(plan);
\draw[uaOrangeArrow] (plan)--(run);
\end{tikzpicture}
\end{center}
\end{frame}

\begin{frame}{Cómo leer un plan: de abajo hacia arriba}
\UASlideTag{NÚCLEO}{UAIngNavy}
\begin{center}
\begin{tikzpicture}[node distance=0.38cm]
\node[uaOrange,minimum width=6.2cm] (limit) {LIMIT: entrega 20 filas};
\node[uaSoft,below=of limit,minimum width=6.2cm] (sort) {SORT: ordena por fecha descendente};
\node[uaSoft,below=of sort,minimum width=6.2cm] (filter) {FILTER: conserva estudiante y estado};
\node[uaSoft,below=of filter,minimum width=6.2cm] (scan) {SEQ SCAN: visita la tabla};
\draw[uaArrow] (scan)--(filter);
\draw[uaArrow] (filter)--(sort);
\draw[uaOrangeArrow] (sort)--(limit);
\node[uaSoft,text width=3.7cm,align=left,right=0.9cm of filter] (read) {Empiece por el nodo inferior: allí aparece el trabajo físico inicial. Luego suba y pregunte cuánto trabajo descarta o transforma cada nodo.};
\draw[uaArrow] (read.west)--(filter.east);
\end{tikzpicture}
\end{center}
\end{frame}

\begin{frame}{Sin índice: mucho trabajo para encontrar 20 filas}
\UASlideTag{NÚCLEO}{UAIngNavy}
\begin{columns}[T]
\column{0.58\textwidth}
\begin{beamercolorbox}[wd=\linewidth,sep=0.18cm,rounded=true,shadow=false]{block body}
\ttfamily\scriptsize
Limit  (actual time=12.55..12.58 rows=20)\\
\quad -> Sort  (rows=20)\\
\qquad Sort Key: confirmed\_at DESC\\
\qquad -> Seq Scan on purchases\\
\qquad\quad Filter: student\_id=8421\\
\qquad\quad Rows Removed by Filter: 9980\\
\qquad\quad Buffers: shared hit=840\\
Execution Time: 12.84 ms
\end{beamercolorbox}
\column{0.38\textwidth}
\begin{uawarn}
\textbf{Tres evidencias}
\begin{itemize}
\item Seq Scan.
\item 9980 filas descartadas.
\item 840 buffers tocados.
\end{itemize}
\end{uawarn}
\end{columns}
\begin{block}{Interpretación}
El motor recorre, filtra y ordena un conjunto mucho mayor que las 20 filas finales.
\end{block}
\end{frame}

\begin{frame}{Con índice compuesto: el plan hace menos trabajo}
\UASlideTag{NÚCLEO}{UAIngNavy}
\scriptsize
\begin{columns}[T]
\column{0.60\textwidth}
\begin{beamercolorbox}[wd=\linewidth,sep=0.15cm,rounded=true,shadow=false]{block body}
\ttfamily\scriptsize
Limit  (actual time=0.20..0.24 rows=20)\\
\quad -> Index Scan using\\
\qquad idx\_purchase\_student\_status\_date\\
\qquad on purchases\\
\qquad Index Cond: student\_id=8421\\
\qquad\quad AND status='CONFIRMED'\\
\qquad Buffers: shared hit=5\\
Execution Time: 0.27 ms
\end{beamercolorbox}
\vspace{0.15cm}
\begin{block}{DDL utilizado}
\texttt{CREATE INDEX ... ON purchases(student\_id, status, confirmed\_at DESC);}
\end{block}
\column{0.36\textwidth}
\begin{uakeybox}
\textbf{Qué cambió}
\begin{itemize}
\item Se tocaron 5 buffers.
\item Se obtuvieron 20 filas útiles.
\item No hubo ordenamiento separado.
\item \texttt{LIMIT} detuvo el recorrido temprano.
\end{itemize}
\end{uakeybox}
\end{columns}
\end{frame}

\begin{frame}{Comparar planes sin convertir los números en una promesa universal}
\UASlideTag{NÚCLEO}{UAIngNavy}
\begin{center}
\begin{tabular}{lcc}
\toprule
Evidencia & Sin índice & Con índice \\
\midrule
Acceso & Seq Scan & Index Scan \\
Filas descartadas & 9980 & 0 en el ejemplo \\
Buffers & 840 & 5 \\
Sort explícito & Sí & No \\
Tiempo ilustrativo & 12.84 ms & 0.27 ms \\
\bottomrule
\end{tabular}
\end{center}
\begin{alertblock}{Cuidado}
Los valores dependen de datos, estadísticas, caché, selectividad y hardware. La lección es aprender a leer evidencia, no memorizar una cifra.
\end{alertblock}
\end{frame}

\begin{frame}{El costo del índice aparece en cada escritura}
\UASlideTag{NÚCLEO}{UAIngNavy}
\begin{center}
\begin{tikzpicture}[node distance=0.5cm and 0.58cm]
\node[uaOrange,minimum width=2.0cm] (ins) {\texttt{INSERT}};
\node[uaSoft,right=of ins,minimum width=2.0cm] (heap) {tabla};
\node[uaSoft,right=of heap,minimum width=2.0cm] (i1) {índice 1};
\node[uaSoft,right=of i1,minimum width=2.0cm] (i2) {índice 2};
\node[uaSoft,right=of i2,minimum width=2.0cm] (i3) {índice 3};
\draw[uaOrangeArrow] (ins)--(heap);
\draw[uaArrow] (heap)--(i1);
\draw[uaArrow] (i1)--(i2);
\draw[uaArrow] (i2)--(i3);
\node[uaTiny,below=0.35cm of heap]{más páginas};
\node[uaTiny,below=0.35cm of i1]{más WAL};
\node[uaTiny,below=0.35cm of i2]{más caché};
\node[uaTiny,below=0.35cm of i3]{más mantenimiento};
\end{tikzpicture}
\end{center}
\begin{block}{Regla de ingeniería}
No se agrega un índice porque ``parece útil''. Se formula una hipótesis, se observa el plan y se mide la carga de trabajo real.
\end{block}
\end{frame}

% =============================================================================
\UASectionSlide{LSM-tree: escrituras rápidas,\\organización posterior}

\begin{frame}{Qué es una LSM-tree y por qué aparece}
\UASlideTag{NÚCLEO}{UAIngNavy}
\small
\begin{columns}[T]
\column{0.54\textwidth}
\begin{block}{Definición intuitiva}
Una LSM-tree recibe escrituras en memoria y luego las vuelca a archivos ordenados e inmutables. Como esos archivos se acumulan, el motor debe fusionarlos periódicamente.
\end{block}
\begin{uakeybox}
No elimina el costo de escritura: desplaza gran parte de la entrada/salida aleatoria fuera del camino crítico y la paga mediante mantenimiento posterior.
\end{uakeybox}
\column{0.42\textwidth}
\begin{center}
\begin{tikzpicture}[node distance=0.22cm]
\node[uaOrange,minimum width=3.15cm] (wal) {WAL};
\node[uaSoft,minimum width=3.15cm,below=of wal] (mem) {tabla en memoria};
\node[uaProcess,minimum width=2.20cm,below=of mem] (flush) {volcado (flush)};
\node[uaSoft,minimum width=3.15cm,below=of flush] (files) {SSTables en L0};
\node[uaProcess,minimum width=2.20cm,below=of files] (compact) {compactación};
\node[uaSoft,minimum width=3.15cm,below=of compact] (levels) {niveles L1, L2, ...};
\draw[uaOrangeArrow] (wal)--(mem);
\draw[uaArrow] (mem)--(flush);
\draw[uaOrangeArrow] (flush)--(files);
\draw[uaArrow] (files)--(compact);
\draw[uaOrangeArrow] (compact)--(levels);
\end{tikzpicture}
\end{center}
\end{columns}
\end{frame}

\begin{frame}{Vocabulario LSM antes de hablar de compactación}
\UASlideTag{NÚCLEO}{UAIngNavy}
\small
\begin{center}
\begin{tabular}{p{0.22\textwidth}p{0.64\textwidth}}
\toprule
Término & Significado \\
\midrule
Memtable & Tabla ordenada en memoria que recibe las escrituras recientes. \\
Flush & Volcado de una memtable inmutable a un archivo persistente. \\
SSTable & \textit{Sorted String Table}: archivo ordenado e inmutable de pares clave--valor. \\
Run & Secuencia o archivo ya ordenado que puede participar en una fusión. \\
Nivel L0, L1, L2 & Conjuntos sucesivos de SSTables organizados por tamaño y solapamiento. \\
Marca de borrado (tombstone) & Registro que representa una eliminación y debe conservarse hasta que sea seguro olvidar versiones anteriores. \\
\bottomrule
\end{tabular}
\end{center}
\end{frame}

\begin{frame}{Vocabulario LSM: mantenimiento y protección}
\UASlideTag{NÚCLEO}{UAIngNavy}
\small
\begin{center}
\begin{tabular}{p{0.22\textwidth}p{0.64\textwidth}}
\toprule
Término & Significado \\
\midrule
Merge ordenado & Fusión que recorre runs ya ordenados, elige la menor clave disponible y produce otra secuencia ordenada sin ordenar desde cero. \\
Compactación & Merge ordenado que además reduce solapamiento y elimina versiones obsoletas cuando es seguro. \\
Trabajo pendiente (backlog) & Cantidad de bytes o archivos que todavía deben ser compactados. \\
Pausa de escritura (write stall) & Desaceleración o detención temporal de escrituras que protege al motor cuando el trabajo pendiente amenaza la estabilidad. \\
Amplificación & Trabajo físico adicional respecto del dato lógico: puede medirse en escrituras, lecturas o espacio. \\
\bottomrule
\end{tabular}
\end{center}
\end{frame}

\begin{frame}{Caso contrastante: telemetría del campus}
\UASlideTag{NÚCLEO}{UAIngNavy}
\scriptsize
\begin{uakeybox}
Sensores de aulas, energía y ocupación producen 80.000 eventos por segundo. Las consultas recorren un dispositivo y una ventana temporal.
\end{uakeybox}
\begin{center}
\begin{tikzpicture}[node distance=0.30cm and 0.55cm]
\node[uaSoft,text width=2.35cm] (s1) {sensor de aula};
\node[uaSoft,text width=2.35cm,below=of s1] (s2) {medidor de energía};
\node[uaSoft,text width=2.35cm,below=of s2] (s3) {sensor climático};
\node[uaOrange,text width=2.65cm,right=0.80cm of s2] (in) {canal de ingestión\\80.000 eventos/s};
\node[uaSoft,text width=3.0cm,right=0.80cm of in] (key) {clave ordenada\\(dispositivo, tiempo)};
\node[uaSoft,text width=2.65cm,right=0.80cm of key] (store) {motor LSM\\retención temporal};
\draw[uaArrow] (s1.east)--(in.west);
\draw[uaOrangeArrow] (s2.east)--(in.west);
\draw[uaArrow] (s3.east)--(in.west);
\draw[uaOrangeArrow] (in)--(key);
\draw[uaArrow] (key)--(store);
\end{tikzpicture}
\end{center}
\begin{columns}[T]
\column{0.48\textwidth}
\begin{itemize}
\item Tasa de escritura alta y sostenida.
\item Claves ordenables por dispositivo y tiempo.
\end{itemize}
\column{0.48\textwidth}
\begin{itemize}
\item Retención y borrado por antigüedad.
\item Búsquedas puntuales y rangos temporales.
\end{itemize}
\end{columns}
\end{frame}

\begin{frame}{Ruta de escritura LSM: del write a una SSTable L0}
\UASlideTag{NÚCLEO}{UAIngNavy}
\small
\begin{center}
\begin{tikzpicture}[node distance=0.28cm and 0.30cm]
\node[uaOrange,text width=1.55cm] (w) {escritura\\$(k,v)$};
\node[uaSoft,text width=1.35cm,right=of w] (wal) {WAL};
\node[uaSoft,text width=1.65cm,right=of wal] (mem) {memtable};
\node[uaSoft,text width=2.05cm,right=of mem] (imm) {memtable\\inmutable};
\node[uaProcess,text width=1.45cm,right=of imm] (flush) {volcado\\(flush)};
\node[uaSoft,text width=1.65cm,right=of flush] (l0) {SSTable\\nivel L0};
\draw[uaOrangeArrow] (w)--(wal);
\draw[uaArrow] (wal)--(mem);
\draw[uaArrow] (mem)--(imm);
\draw[uaOrangeArrow] (imm)--(flush);
\draw[uaOrangeArrow] (flush)--(l0);
\end{tikzpicture}
\end{center}
\begin{block}{Qué ocurre en el camino crítico}
La escritura se registra en el WAL, se incorpora a una estructura ordenada en memoria y, más tarde, una memtable inmutable se vuelca como una SSTable de nivel L0.
\end{block}
\begin{alertblock}{Qué queda para después}
L0 puede acumular archivos solapados. La compactación posterior reorganiza esos archivos y los desplaza hacia niveles más estables.
\end{alertblock}
\end{frame}

\begin{frame}{Compactación: merge ordenado y eliminación segura}
\UASlideTag{NÚCLEO}{UAIngNavy}
\small
\begin{center}
\begin{tikzpicture}[x=1cm,y=1cm]
\node[uaSoft,text width=2.15cm] (a) at (-4.70,1.10) {SSTable A\\a, c1, e, g};
\node[uaSoft,text width=2.15cm] (b) at (-4.70,0.00) {SSTable B\\b, c2, f, h};
\node[uaSoft,text width=2.15cm] (t) at (-4.70,-1.10) {marca de borrado\\eliminar c};
\node[uaOrange,text width=3.25cm] (m) at (0.00,0.00) {merge ordenado\\compara claves\\conserva la versión vigente};
\node[uaSoft,text width=2.75cm] (out) at (4.45,0.00) {SSTable resultante\\a, b, e, f, g, h};
\coordinate (mA) at ($(m.west)+(0,0.65)$);
\coordinate (mB) at (m.west);
\coordinate (mT) at ($(m.west)+(0,-0.65)$);
\draw[uaArrow] (a.east)--++(0.45,0)|-(mA);
\draw[uaArrow] (b.east)--(mB);
\draw[uaArrow] (t.east)--++(0.45,0)|-(mT);
\draw[uaOrangeArrow] (m)--(out);
\end{tikzpicture}
\end{center}
\begin{block}{Qué se paga}
La compactación lee y reescribe datos, utiliza espacio temporal, compite por E/S y puede aumentar el percentil 99 de latencia.
\end{block}
\end{frame}

\begin{frame}{Filtro de Bloom: decidir qué SSTables se pueden omitir}
\UASlideTag{NÚCLEO}{UAIngNavy}
\scriptsize
\begin{uainfo}
Un \textbf{filtro de Bloom} es un resumen probabilístico y compacto de las claves de una SSTable. Para una búsqueda de clave exacta no devuelve el valor: responde \textbf{definitivamente ausente} o \textbf{posiblemente presente}.
\end{uainfo}
\begin{center}
\begin{tikzpicture}[node distance=0.34cm and 0.40cm]
\node[uaSoft,text width=2.45cm] (sst) {SSTable A\\claves: a, c, e, g};
\node[uaOrange,text width=2.45cm,right=of sst] (bf) {Filtro de Bloom\\resumen de claves};
\node[uaSoft,text width=2.15cm,right=of bf,yshift=0.72cm] (qno) {buscar z};
\node[uaSoft,text width=2.15cm,right=of bf,yshift=-0.72cm] (qmaybe) {buscar c};
\node[uaOrange,text width=2.45cm,right=of qno] (no) {algún bit es 0\\definitivamente ausente};
\node[uaSoft,text width=2.45cm,right=of qmaybe] (maybe) {todos los bits son 1\\posiblemente presente};
\draw[uaOrangeArrow] (sst)--(bf);
\draw[uaArrow] (bf.east)--(qno.west);
\draw[uaArrow] (bf.east)--(qmaybe.west);
\draw[uaOrangeArrow] (qno)--(no);
\draw[uaArrow] (qmaybe)--(maybe);
\end{tikzpicture}
\end{center}
\begin{itemize}
\item \textbf{Definitivamente ausente:} el motor omite esa SSTable para la búsqueda. Esto reemplaza la expresión ambigua ``archivo imposible''.
\item \textbf{Posiblemente presente:} el motor consulta el índice y, si corresponde, los bloques de datos; puede ser un falso positivo.
\item Un filtro bien construido no debe producir falsos negativos: no puede descartar un archivo que sí contiene la clave.
\end{itemize}
\end{frame}

\begin{frame}{Ruta de lectura LSM: descartar antes de leer el dispositivo}
\UASlideTag{NÚCLEO}{UAIngNavy}
\scriptsize
\begin{center}
\resizebox{0.96\textwidth}{!}{%
\begin{tikzpicture}[node distance=0.30cm and 0.34cm]
\node[uaOrange,text width=1.70cm] (get) {buscar clave $k$};
\node[uaSoft,text width=1.70cm,right=of get] (mem) {revisar memtables};
\node[uaSoft,text width=2.05cm,right=of mem] (cand) {SSTable candidata};
\node[diamond,draw=UAIngNavy,thick,fill=UAIngPaper,aspect=2.15,align=center,font=\tiny,
      inner xsep=1.5pt,inner ysep=1pt,right=0.42cm of cand] (bf) {¿Bloom dice\\``tal vez''?};
\node[uaSoft,text width=2.05cm,right=0.45cm of bf,yshift=0.72cm] (skip) {No\\omitir SSTable};
\node[uaOrange,text width=2.05cm,right=0.45cm of bf,yshift=-0.72cm] (idx) {Sí\\consultar índice};
\node[uaSoft,text width=2.15cm,right=0.35cm of idx] (data) {caché o\\bloque en dispositivo};
\draw[uaOrangeArrow] (get)--(mem);
\draw[uaArrow] (mem)--(cand);
\draw[uaArrow] (cand)--(bf);
\draw[uaArrow] (bf.east)--(skip.west);
\draw[uaOrangeArrow] (bf.east)--(idx.west);
\draw[uaOrangeArrow] (idx)--(data);
\end{tikzpicture}}%
\end{center}
\vspace{0.08cm}
\begin{columns}[T]
\column{0.48\textwidth}
\begin{block}{Orden de búsqueda}
Primero se revisa lo más reciente. Para cada SSTable candidata, Bloom puede evitar una lectura física.
\end{block}
\column{0.48\textwidth}
\begin{block}{Resultado correcto}
Si aparecen varias versiones, el motor elige la más reciente que sea visible y respeta las marcas de borrado.
\end{block}
\end{columns}
\begin{itemize}
\item Un falso positivo agrega trabajo; un falso negativo haría perder un dato existente y no es admisible.
\item La caché de bloques evita acceder al dispositivo cuando el bloque requerido ya está en memoria.
\end{itemize}
\end{frame}

\begin{frame}{Panel de compactación I: qué señales observar}
\UASlideTag{NÚCLEO}{UAIngNavy}
\scriptsize
\renewcommand{\arraystretch}{0.95}
\begin{center}
\begin{tabular}{p{0.34\textwidth}p{0.17\textwidth}p{0.36\textwidth}}
\toprule
Señal & Valor ilustrativo & Qué informa \\
\midrule
Ingreso lógico & 90 MB/s & Trabajo nuevo recibido por el motor. \\
Capacidad de compactación & 60 MB/s & Ritmo de reorganización de archivos. \\
Déficit sostenido & 30 MB/s & Trabajo que se acumula por segundo. \\
Trabajo pendiente & 18 GB & Datos todavía sin compactar. \\
Archivos en L0 & 28 de 36 & Cercanía al umbral de protección. \\
Percentil 99 de escritura & En aumento & Impacto sobre las operaciones más lentas. \\
\bottomrule
\end{tabular}
\end{center}
\vspace{0.08cm}
\textbf{Lectura conjunta.} Relacione ingreso, capacidad de mantenimiento, crecimiento del trabajo pendiente y latencia; ninguna cifra aislada basta para diagnosticar.
\end{frame}

\begin{frame}{Panel de compactación II: leer la evolución}
\UASlideTag{NÚCLEO}{UAIngNavy}
\small
\begin{center}
\begin{tikzpicture}[x=0.72cm,y=0.48cm]
\node[anchor=west,font=\small\bfseries,text=UAIngNavy] at (0,6.00){Trabajo pendiente de compactación durante 10 minutos};
\draw[->,thick,UAIngNavy] (0.8,0.65)--(11.1,0.65) node[right,uaTiny]{minutos};
\draw[->,thick,UAIngNavy] (0.8,0.65)--(0.8,4.75) node[above,uaTiny]{GB};
\foreach \x/\h in {1.15/0.45,2.20/0.82,3.25/1.22,4.30/1.68,5.35/2.15,6.40/2.65,7.45/3.15,8.50/3.62}{
  \fill[UAIngOrange] (\x,0.65) rectangle +(0.58,\h);
}
\draw[UAIngNavy,dashed,thick] (0.8,3.20)--(9.35,3.20);
\draw[UAIngOrange,dashed,thick] (0.8,4.30)--(9.35,4.30);
\node[uaTiny,anchor=west] at (9.55,3.20){umbral de desaceleración};
\node[uaTiny,anchor=west,text=UAIngOrange] at (9.55,4.30){umbral de detención};
\node[uaTiny,anchor=west] at (1.15,0.22){0};
\node[uaTiny,anchor=west] at (8.50,0.22){10};
\end{tikzpicture}
\end{center}
\begin{block}{Cadena causal}
Ingreso mayor que capacidad de compactación $\Rightarrow$ trabajo pendiente creciente $\Rightarrow$ más archivos en L0 $\Rightarrow$ mayor trabajo de lectura y latencia $\Rightarrow$ protección mediante una pausa de escritura.
\end{block}
\end{frame}

\begin{frame}{Pausa de escritura (write stall): proteger la estabilidad}
\UASlideTag{NÚCLEO}{UAIngNavy}
\small
\begin{center}
\begin{tikzpicture}[node distance=0.30cm and 0.38cm]
\node[uaOrange,text width=2.25cm] (in) {Ingreso solicitado\\90 MB/s};
\node[uaSoft,text width=2.25cm,right=of in] (cap) {Capacidad de compactación\\60 MB/s};
\node[uaOrange,text width=2.05cm,right=of cap] (def) {Déficit sostenido\\30 MB/s};
\draw[uaOrangeArrow] (in)--(cap);
\draw[uaOrangeArrow] (cap)--(def);

\node[uaSoft,text width=2.05cm,below=0.70cm of def] (q) {Trabajo pendiente\\crece};
\node[uaSoft,text width=2.05cm,left=0.38cm of q] (th) {Se supera\\un umbral};
\node[uaOrange,text width=2.05cm,left=0.38cm of th] (stall) {Pausa de escritura\\reduce admisión};
\node[uaSoft,text width=2.45cm,left=0.38cm of stall] (stable) {Ingreso efectivo $\leq$ capacidad\\el backlog se estabiliza};
\draw[uaArrow] (def)--(q);
\draw[uaArrow] (q)--(th);
\draw[uaOrangeArrow] (th)--(stall);
\draw[uaOrangeArrow] (stall)--(stable);
\end{tikzpicture}
\end{center}
\begin{alertblock}{Interpretación correcta}
La pausa no crea capacidad: evita que una cola interna crezca sin límite. Si aparece con frecuencia o dura demasiado, revela una carga sostenida superior a la capacidad o una política de compactación inadecuada.
\end{alertblock}
\end{frame}

\begin{frame}{Leveled y tiered: dos formas de distribuir el costo}
\UASlideTag{NÚCLEO}{UAIngNavy}
\small
\begin{center}
\begin{tabular}{p{0.23\textwidth}p{0.32\textwidth}p{0.32\textwidth}}
\toprule
 & Por niveles (leveled) & Por grupos (tiered / universal) \\
\midrule
Lecturas & Menos archivos solapados & Más archivos potenciales \\
Escrituras & Más reescritura en merges & Menos reescritura temprana \\
Espacio & Mejor control & Más espacio temporal \\
Carga típica & Búsquedas puntuales sensibles & Ingestión intensa \\
\bottomrule
\end{tabular}
\end{center}
\begin{block}{Decisión}
La política depende del patrón de acceso, del hardware, de la retención, de las amplificaciones y del objetivo de nivel de servicio (SLO); no del nombre del producto.
\end{block}
\end{frame}

% =============================================================================
\UASectionSlide{Integración: decidir con evidencia}

\begin{frame}{B-tree o LSM: una orientación, no una ley universal}
\UASlideTag{NÚCLEO}{UAIngNavy}
\scriptsize
\begin{center}
\begin{tabular}{p{0.17\textwidth}p{0.35\textwidth}p{0.35\textwidth}}
\toprule
Criterio & Árbol B y páginas & LSM y archivos ordenados \\
\midrule
Camino de escritura & Actualiza páginas e índices & WAL + memtable + flush + compactación \\
Búsqueda puntual & Pocas páginas si el índice ayuda & Puede revisar varios archivos; filtros y caché ayudan \\
Rangos & Natural por el orden del árbol & Eficiente si la compactación mantiene pocos archivos \\
Mantenimiento & Divisiones, VACUUM, bloat y reconstrucción & Tombstones, compactación y write stalls \\
Carga típica & OLTP (procesamiento transaccional en línea) o predominio de lecturas & Ingestión sostenida y predominio de escrituras \\
Pregunta crítica & ¿Qué consulta justifica el índice? & ¿Puede la compactación sostener el ritmo? \\
\bottomrule
\end{tabular}
\end{center}
\end{frame}

\begin{frame}{Dos casos cercanos, dos hipótesis defendibles}
\UASlideTag{NÚCLEO}{UAIngNavy}
\scriptsize
\begin{columns}[T]
\column{0.48\textwidth}
\begin{uainfo}
\textbf{CampusShop}
\begin{itemize}
\item Transacciones cortas.
\item Inventario e idempotencia.
\item Consultas por estudiante y fecha.
\item Percentil 99 predecible.
\end{itemize}
\textbf{Hipótesis:} WAL + MVCC + árbol B.
\end{uainfo}
\column{0.48\textwidth}
\begin{uainfo}
\textbf{Telemetría del campus}
\begin{itemize}
\item Ingestión masiva y secuencial.
\item Rangos por dispositivo y tiempo.
\item Retención temporal.
\item Vistas con frescura eventual.
\end{itemize}
\textbf{Hipótesis:} WAL + LSM + compactación temporal.
\end{uainfo}
\end{columns}
\begin{block}{Todavía falta evidencia}
Medir percentiles 95 y 99, throughput sostenible, objetivo de punto de recuperación (RPO), objetivo de tiempo de recuperación (RTO), amplificación, backlog, tasa de aciertos de caché y espacio.
\end{block}
\end{frame}

\begin{frame}{Actividad de 12 minutos: problema, evidencia, decisión}
\UASlideTag{ACTIVIDAD}{UAIngOrange}
\small
\begin{columns}[T]
\column{0.48\textwidth}
\begin{uainfo}
\textbf{Equipo A --- CampusShop}
\begin{itemize}
\item Explique el plan sin índice.
\item Defienda o rechace el índice compuesto.
\item Proponga dos métricas de durabilidad y consulta.
\end{itemize}
\end{uainfo}
\column{0.48\textwidth}
\begin{uainfo}
\textbf{Equipo B --- Telemetría}
\begin{itemize}
\item Interprete backlog, amplificación de escritura y write stall.
\item Proponga una hipótesis de causa.
\item Diga qué medir antes de ajustar la configuración.
\end{itemize}
\end{uainfo}
\end{columns}
\begin{alertblock}{Entrega oral}
Cada equipo dispone de 90 segundos para presentar el problema, la evidencia, la decisión y su trade-off.
\end{alertblock}
\end{frame}

\begin{frame}{Volvemos al diagnóstico inicial}
\UASlideTag{ACTIVIDAD}{UAIngOrange}
\begin{enumerate}
\item \textbf{Falso:} el COMMIT puede ser durable aunque la página siga atrasada.
\item \textbf{Falso:} copiar en memoria principal no equivale a persistir.
\item \textbf{Falso:} MVCC reduce contención, pero no elimina todas las anomalías.
\item \textbf{Falso:} un índice puede no utilizarse y siempre tiene costo de mantenimiento.
\item \textbf{Falso:} una LSM desplaza trabajo hacia el flush y la compactación.
\end{enumerate}
\begin{block}{Patrón común}
Ningún mecanismo elimina el costo: cambia \emph{cuándo}, \emph{dónde} y \emph{cómo} se paga.
\end{block}
\end{frame}

\begin{frame}{Método reusable para decidir almacenamiento}
\UASlideTag{NÚCLEO}{UAIngNavy}
\begin{center}
\begin{tikzpicture}[node distance=0.4cm]
\node[uaOrange,text width=2.0cm] (a) {1. Carga};
\node[uaSoft,text width=2.0cm,right=of a] (b) {2. Garantía};
\node[uaSoft,text width=2.0cm,right=of b] (c) {3. Mecanismo};
\node[uaSoft,text width=2.0cm,right=of c] (d) {4. Costo};
\node[uaSoft,text width=2.0cm,right=of d] (e) {5. Evidencia};
\draw[uaOrangeArrow] (a)--(b);
\draw[uaArrow] (b)--(c);
\draw[uaArrow] (c)--(d);
\draw[uaArrow] (d)--(e);
\draw[uaOrangeArrow] (e) to[bend left=27] node[above,uaTiny]{revisar} (a);
\end{tikzpicture}
\end{center}
\begin{itemize}
\item Describa proporciones de lectura y escritura, tamaño, sesgo, crecimiento y consultas.
\item Declare qué falla debe tolerarse y qué significa durable en ese modelo.
\item Elija un mecanismo inicial, no un producto por reputación.
\item Identifique el costo que se desplaza a otra etapa.
\item Diseñe un experimento capaz de contradecir su elección.
\end{itemize}
\end{frame}

\begin{frame}{Cierre: de la intuición a la evidencia}
\UASlideTag{NÚCLEO}{UAIngNavy}
\begin{center}
\begin{tikzpicture}[node distance=0.38cm]
\node[uaSoft,text width=2.05cm] (g) {promesa de negocio};
\node[uaSoft,text width=2.05cm,right=of g] (wal) {WAL\\permite recuperar};
\node[uaSoft,text width=2.05cm,right=of wal] (mvcc) {MVCC\\mantiene vistas};
\node[uaSoft,text width=2.05cm,right=of mvcc] (idx) {índice + EXPLAIN\\reducen trabajo};
\node[uaOrange,text width=2.05cm,right=of idx] (lsm) {compactación + pausa\\protegen capacidad};
\draw[uaArrow] (g)--(wal);
\draw[uaArrow] (wal)--(mvcc);
\draw[uaArrow] (mvcc)--(idx);
\draw[uaOrangeArrow] (idx)--(lsm);
\end{tikzpicture}
\end{center}
\begin{block}{Idea final}
El almacenamiento no es una caja debajo del sistema. Es la capa donde las promesas lógicas se convierten en bytes, esperas, recuperación y riesgos medibles.
\end{block}
\end{frame}

\begin{frame}{Exit ticket}
\UASlideTag{ACTIVIDAD}{UAIngOrange}
Responda en tres líneas por punto:
\begin{enumerate}
\item ¿Qué diferencia hay entre \emph{aceptado}, \emph{escrito} y \emph{durable}?
\item ¿Qué costo oculto introduce MVCC y cómo lo observaría?
\item ¿Qué diferencia hay entre \texttt{EXPLAIN} y \texttt{EXPLAIN ANALYZE}?
\item ¿Qué señales indican que una LSM acumula trabajo pendiente de compactación?
\end{enumerate}
\begin{alertblock}{Puente a la clase 7}
Hoy vimos cómo un nodo conserva y organiza su historia. La próxima clase estudia cómo un log se convierte en una interfaz distribuida entre productores y consumidores.
\end{alertblock}
\end{frame}

\end{document}
